\chapter{User Interface Design}
\ChapterLabel{UserInterfaceDesign}
\section{Chapter Responsibility}
\label{Section:UserInterfaceDesignResponsibility}
Define how people interact with the system and understand its state.

\noindent\textbf{Expected result:} An interaction design that supports stakeholder tasks and acceptance.

\section{SP800Project Status}
\label{Section:SP800StatusUserInterfaceDesign}
Project-specific work in this chapter is planned. The inherited examples,
identifiers, values, mappings, and statuses below are template reference
material only. They are not approved SP800Project obligations, implemented features,
executed tests, verification evidence, validation, or release authorization.
Dependencies and missing facts are owned in
\hyperref[Section:SP800OpenQuestions]{the project open-question list}.

\section{How to Use This Chapter}
Describe personas, primary screens or workflows, interaction rules, feedback, accessibility, and essential states. Include wireframes or prototypes where they resolve meaningful ambiguity.
\section{Example}
Persona: a busy operations coordinator needs to submit and track requests quickly. The request screen groups required fields, identifies inline validation errors, confirms successful submission, supports keyboard operation, and exposes empty, loading, error, and completed states.
